% (a) Legge di Stevino: colonna di liquido di altezza h e base S.
% (b) Vasi comunicanti di forma diversa: stesso livello, stessa pressione alla stessa quota.
\begin{tikzpicture}[>=Stealth, font=\small]
  % (a) recipiente
  \fill[formula!15] (0,0) rectangle (3,2.8);
  \draw[very thick] (0,3.2) -- (0,0) -- (3,0) -- (3,3.2);
  \draw[formula, thick] (0,2.8) -- (3,2.8);
  % colonna
  \draw[dashed, thick, esempio] (1.1,0.8) rectangle (1.9,2.8);
  \draw[esempio, thick] (1.1,0.8) -- (1.9,0.8);
  \node[left, esempio] at (1.1,0.8) {$S$};
  % frecce
  \draw[->, thick, attenzione] (1.5,3.55) -- (1.5,2.85) node[pos=0, above] {$p_0$};
  \draw[->, thick, attenzione] (1.5,0.1) -- (1.5,0.75);
  \node[attenzione] at (1.85,0.3) {$p$};
  % quota h
  \draw[|<->|] (2.4,0.8) -- (2.4,2.8) node[midway, right] {$h$};
  \node at (1.5,-0.45) {(a)};

  % (b) vasi comunicanti
  \begin{scope}[xshift=4.2cm]
    \fill[formula!15] (0,0) -- (0,2.2) -- (0.8,2.2) -- (0.8,0.5) -- (2.2,0.5) -- (1.757,2.2)
      -- (3.243,2.2) -- (2.8,0.5) -- (4.8,0.5) -- (5.539,2.2) -- (6.039,2.2) -- (5.3,0.5) -- (5.3,0) -- cycle;
    \draw[very thick] (0.8,2.8) -- (0.8,0.5) -- (2.2,0.5) -- (1.6,2.8);
    \draw[very thick] (3.4,2.8) -- (2.8,0.5) -- (4.8,0.5) -- (5.8,2.8);
    \draw[very thick] (0,2.8) -- (0,0) -- (5.3,0) -- (5.3,0.5) -- (6.3,2.8);
    \draw[formula, thick] (0,2.2) -- (0.8,2.2) (1.757,2.2) -- (3.243,2.2) (5.539,2.2) -- (6.039,2.2);
    \draw[dashed, richiamo] (-0.3,2.2) -- (6.6,2.2);
    \draw[dashed, richiamo] (-0.3,1.2) -- (6.6,1.2) node[right, text=black, align=left] {stessa\\quota};
    \foreach \x/\n in {0.4/A, 2.5/B, 5.354/C} {
      \fill (\x,1.2) circle (1.6pt);
      \node[above] at (\x,1.2) {\n};
    }
    \node at (3.0,-0.45) {(b)};
  \end{scope}
\end{tikzpicture}
